\section{Data Stewardship and Compute Workflow}
All data handling activities will respect Wits and hospital governance requirements. Access-controlled mirrors of NIH ChestX-ray14, CheXpert, and any local radiography datasets will be staged on encrypted storage. Metadata versioning will be managed with \texttt{datalad} so that provenance and consent restrictions remain visible throughout the project lifecycle. Preprocessing pipelines will execute on \texttt{bigbatch}, where large memory nodes simplify batched pixel standardisation, report parsing, and anatomical segmentation. Outputs will be validated against held-out clinical samples before being promoted to the research data lake. Feature statistics, quality reports, and lineage diagrams will be archived alongside the processed datasets to support reproducibility and facilitate future audits.

To optimise utilisation of the HPC environment, I will orchestrate experiments through a queue-aware scheduler that distributes workloads across \texttt{bigbatch} for data-intensive tasks, \texttt{stampede} for mid-scale hyper-parameter sweeps, and \texttt{biggpu} for large-scale diffusion training. Each submission script will capture environment hashes, resource requests, and runtime metrics, enabling post-hoc performance analysis. Regular clean-up cycles and usage dashboards will prevent quota overruns and ensure that compute-intensive campaigns remain aligned with project priorities.